% GENERATED FILE. Edit canonical lessons, then run npm run generate:books.
% canonical-source-hash: fnv1a64:00f69b5642d0cf2f
% canonical-lessons: TA-C379-tannarvalar, TA-C379-payanali, TA-C379-niti-funds, TA-C379-tirattu, TA-C379-nalan

\chapter{Helping Others}
\label{ch:ta-helping-others}
\hlchaptermodality{379}

\begin{chapteropening}
\textbf{By the end of this chapter:} \emph{I can talk about volunteers, funds and welfare.}

Complete the last lesson of chapter 379: I can talk about volunteers, funds and welfare.
\end{chapteropening}

\section[\texorpdfstring{\ta{தன்னார்வலர்} (taṉṉārvalar)}{தன்னார்வலர் (taṉṉārvalar)}]{\ta{தன்னார்வலர்} (taṉṉārvalar) — a volunteer}

\label{lesson:TA-C379-tannarvalar}



\begin{quote}
Before the new one: say the Tamil for a leader, then the Tamil for a treasurer.
\end{quote}

\subsection*{You'll want to know: \ta{தன்னார்வலர்}}
\textbf{\ta{தன்னார்வலர்}} — \emph{taṉṉārvalar} — \textquotedblleft{}a volunteer\textquotedblright{}.

\begin{quote}
\ta{வெள்ளத்தின்போது} \ta{பல} \ta{தன்னார்வலர்கள்} \ta{உதவினார்கள்}.
\end{quote}

\begin{quote}
\emph{veḷḷattiṉpōtu pala taṉṉārvalarkaḷ utaviṉārkaḷ.}
\end{quote}

\begin{quote}
Many volunteers helped during the flood.
\end{quote}

\ta{தன்}, self + \ta{ஆர்வலர்}, an eager person: one who comes of their own will.

\begin{cousinweb}
\textbf{\ta{தன்னார்வலர்}} (\emph{taṉṉārvalar}) is a modern native coinage: \emph{taṉ}, \textquotedblleft{}oneself\textquotedblright{} (DEDR 3196), and \emph{ārvalar}, \textquotedblleft{}the eager\textquotedblright{}, from \emph{ārvam}, \textquotedblleft{}desire, devotion\textquotedblright{} (DEDR 381).
\end{cousinweb}

\begin{grammarlens}[title={What you've built}]
One more word for giving reasons and opinions.
\end{grammarlens}

\subsection*{Your turn}
\begin{itemize}
\raggedright
  \item \emph{Say it:} \emph{taṉṉārvalar}
  \item \emph{Say it:} \emph{taṉṉārvalar}, once more
  \item \emph{Say it:} \emph{veḷḷattiṉpōtu pala taṉṉārvalarkaḷ utaviṉārkaḷ.} — Many volunteers helped during the flood.
  \item \emph{Recall:} say the Tamil for a leader, then the Tamil for a treasurer, then say \emph{taṉṉārvalar} again
\end{itemize}

\begin{tcolorbox}[breakable,colback=teal!4,colframe=teal!35!black,title={Before you move on}]
What does \ta{தன்னார்வலர்} mean? (\textquotedblleft{}A volunteer\textquotedblright{}.) Say it once more.
\end{tcolorbox}
\section[\texorpdfstring{\ta{பயனாளி} (payaṉāḷi)}{பயனாளி (payaṉāḷi)}]{\ta{பயனாளி} (payaṉāḷi) — a beneficiary}

\label{lesson:TA-C379-payanali}



\begin{quote}
Before the new one: say the Tamil for a treasurer, then the Tamil for a volunteer.
\end{quote}

\subsection*{You'll want to know: \ta{பயனாளி}}
\textbf{\ta{பயனாளி}} — \emph{payaṉāḷi} — \textquotedblleft{}a beneficiary\textquotedblright{}.

\begin{quote}
\ta{இந்தத்} \ta{திட்டத்தின்} \ta{பயனாளிகள்} \ta{யார்}?
\end{quote}

\begin{quote}
\emph{intat tiṭṭattiṉ payaṉāḷikaḷ yār?}
\end{quote}

\begin{quote}
Who are the beneficiaries of this scheme?
\end{quote}

Plural: \ta{பயனாளிகள்}, \emph{payaṉāḷikaḷ}.

\ta{பயன்}, benefit + -\ta{ஆளி}, the one who holds.

\begin{cousinweb}
\textbf{\ta{பயனாளி}} (\emph{payaṉāḷi}) is built from \emph{payaṉ}, \textquotedblleft{}use, benefit\textquotedblright{}, and native \emph{āḷi}, \textquotedblleft{}one who holds\textquotedblright{} (DEDR 399), as in \ta{முதலாளி}.
\end{cousinweb}

\begin{grammarlens}[title={What you've built}]
One more word for giving reasons and opinions.
\end{grammarlens}

\subsection*{Your turn}
\begin{itemize}
\raggedright
  \item \emph{Say it:} \emph{payaṉāḷi}
  \item \emph{Say it:} \emph{payaṉāḷi}, once more
  \item \emph{Say it:} \emph{intat tiṭṭattiṉ payaṉāḷikaḷ yār?} — Who are the beneficiaries of this scheme?
  \item \emph{Recall:} say the Tamil for a treasurer, then the Tamil for a volunteer, then say \emph{payaṉāḷi} again
\end{itemize}

\begin{tcolorbox}[breakable,colback=teal!4,colframe=teal!35!black,title={Before you move on}]
What does \ta{பயனாளி} mean? (\textquotedblleft{}A beneficiary\textquotedblright{}.) Say it once more.
\end{tcolorbox}
\section[\texorpdfstring{\ta{நிதி} (niti)}{நிதி (niti)}]{\ta{நிதி} (niti) — funds; a fund}

\label{lesson:TA-C379-niti-funds}



\begin{quote}
Before the new one: say the Tamil for a volunteer, then the Tamil for a beneficiary.
\end{quote}

\subsection*{You'll want to know: \ta{நிதி}}
\textbf{\ta{நிதி}} — \emph{niti} — \textquotedblleft{}funds; a fund\textquotedblright{}.

\begin{quote}
\ta{பள்ளிக்குப்} \ta{புதிய} \ta{கட்டிடம்} \ta{கட்ட} \ta{நிதி} \ta{தேவை}.
\end{quote}

\begin{quote}
\emph{paḷḷikkup putiya kaṭṭiṭam kaṭṭa niti tēvai.}
\end{quote}

\begin{quote}
Funds are needed to build a new school building.
\end{quote}

\ta{நிதியமைச்சர்}, the finance minister.

\begin{cousinweb}
\textbf{\ta{நிதி}} (\emph{niti}) is borrowed from Sanskrit \emph{nidhi}, \textquotedblleft{}a setting down; hoard, treasure\textquotedblright{} (CDIAL 7207). Tamil writes the dh as \ta{த}.
\end{cousinweb}

\begin{grammarlens}[title={What you've built}]
One more word for giving reasons and opinions.
\end{grammarlens}

\subsection*{Your turn}
\begin{itemize}
\raggedright
  \item \emph{Say it:} \emph{niti}
  \item \emph{Say it:} \emph{niti}, once more
  \item \emph{Say it:} \emph{paḷḷikkup putiya kaṭṭiṭam kaṭṭa niti tēvai.} — Funds are needed to build a new school building.
  \item \emph{Recall:} say the Tamil for a volunteer, then the Tamil for a beneficiary, then say \emph{niti} again
\end{itemize}

\begin{tcolorbox}[breakable,colback=teal!4,colframe=teal!35!black,title={Before you move on}]
What does \ta{நிதி} mean? (\textquotedblleft{}Funds; a fund\textquotedblright{}.) Say it once more.
\end{tcolorbox}
\section[\texorpdfstring{\ta{திரட்டு} (tiraṭṭu)}{திரட்டு (tiraṭṭu)}]{\ta{திரட்டு} (tiraṭṭu) — to muster; to raise (funds)}

\label{lesson:TA-C379-tirattu}



\begin{quote}
Before the new one: say the Tamil for a beneficiary, then the Tamil for funds.
\end{quote}

\subsection*{You'll want to know: \ta{திரட்டு}}
\textbf{\ta{திரட்டு}} — \emph{tiraṭṭu} — \textquotedblleft{}to muster; to raise (funds)\textquotedblright{}.

\begin{quote}
\ta{மாணவர்கள்} \ta{வெள்ள} \ta{நிவாரணத்துக்கு} \ta{நிதி} \ta{திரட்டினார்கள்}.
\end{quote}

\begin{quote}
\emph{māṇavarkaḷ veḷḷa nivāraṇattukku niti tiraṭṭiṉārkaḷ.}
\end{quote}

\begin{quote}
The students raised funds for flood relief.
\end{quote}

Past (I …): \ta{திரட்டினேன்}, \emph{tiraṭṭiṉēṉ}.

\ta{நிதி} \ta{திரட்டு}, to raise funds; \ta{ஆட்களைத்} \ta{திரட்டு}, to gather people.

\begin{cousinweb}
\textbf{\ta{திரட்டு}} (\emph{tiraṭṭu}) is native Dravidian: DEDR 3245 lists Tamil \emph{tiraḷ}, \textquotedblleft{}to assemble, collect in numbers\textquotedblright{}, and its causative \emph{tiraṭṭu}, \textquotedblleft{}to gather\textquotedblright{}.
\end{cousinweb}

\begin{grammarlens}[title={What you've built}]
One more word for giving reasons and opinions.
\end{grammarlens}

\subsection*{Your turn}
\begin{itemize}
\raggedright
  \item \emph{Say it:} \emph{tiraṭṭu}
  \item \emph{Say it:} \emph{tiraṭṭu}, once more
  \item \emph{Say it:} \emph{māṇavarkaḷ veḷḷa nivāraṇattukku niti tiraṭṭiṉārkaḷ.} — The students raised funds for flood relief.
  \item \emph{Recall:} say the Tamil for a beneficiary, then the Tamil for funds, then say \emph{tiraṭṭu} again
\end{itemize}

\begin{tcolorbox}[breakable,colback=teal!4,colframe=teal!35!black,title={Before you move on}]
What does \ta{திரட்டு} mean? (\textquotedblleft{}To muster; to raise (funds)\textquotedblright{}.) Say it once more.
\end{tcolorbox}
\section[\texorpdfstring{\ta{நலன்} (nalaṉ)}{நலன் (nalaṉ)}]{\ta{நலன்} (nalaṉ) — welfare; well-being}

\label{lesson:TA-C379-nalan}



\begin{quote}
Before the new one: say the Tamil for funds, then the Tamil for to muster.
\end{quote}

\subsection*{You'll want to know: \ta{நலன்}}
\textbf{\ta{நலன்}} — \emph{nalaṉ} — \textquotedblleft{}welfare; well-being\textquotedblright{}.

\begin{quote}
\ta{குழந்தைகளின்} \ta{நலனுக்காக} \ta{அவர்} \ta{பல} \ta{வேலைகள்} \ta{செய்தார்}.
\end{quote}

\begin{quote}
\emph{kuḻantaikaḷiṉ nalaṉukkāka avar pala vēlaikaḷ ceytār.}
\end{quote}

\begin{quote}
He did a lot of work for the welfare of children.
\end{quote}

\ta{நலம்}, health, as in \ta{நான்} \ta{நலம்}, is a close cousin.

\begin{cousinweb}
\textbf{\ta{நலன்}} (\emph{nalaṉ}) is native Dravidian, built on \emph{nal}, \textquotedblleft{}good\textquotedblright{} (DEDR 3610), the root of \ta{நலம்}, \textquotedblleft{}well-being\textquotedblright{}, and \ta{நன்மை}, \textquotedblleft{}a benefit\textquotedblright{}.
\end{cousinweb}

\begin{grammarlens}[title={What you've built}]
One more word for giving reasons and opinions.
\end{grammarlens}

\subsection*{Your turn}
\begin{itemize}
\raggedright
  \item \emph{Say it:} \emph{nalaṉ}
  \item \emph{Say it:} \emph{nalaṉ}, once more
  \item \emph{Say it:} \emph{kuḻantaikaḷiṉ nalaṉukkāka avar pala vēlaikaḷ ceytār.} — He did a lot of work for the welfare of children.
  \item \emph{Recall:} say the Tamil for funds, then the Tamil for to muster, then say \emph{nalaṉ} again
\end{itemize}

\begin{tcolorbox}[breakable,colback=teal!4,colframe=teal!35!black,title={Before you move on}]
What does \ta{நலன்} mean? (\textquotedblleft{}Welfare; well-being\textquotedblright{}.) Say it once more.
\end{tcolorbox}
